\section{The Plane and the Sphere}
\label{sec:plane-sphere}


The unit sphere and the plane are two popular surfaces. 
We often think of the sphere and the plane as subsets of $\R^3$,
the sphere is the set
$$\Sp^2=\{(x,y,z)\in \R^3 | x^2+y^2+z^2=1\}$$
and the $xy$-plane is
$$\R^2=\{(x,y,z)\in \R^3 | z=0\}.$$
Is it possible to continuously deform the sphere into the plane?
The answer is no, as we will see.
Stereographic projection is a map between almost all of the sphere
and the plane. Given a point $p$ on the sphere, we draw a straight line
from the north pole $N=(0,0,1)$ through $p$, then find the intersection of
this line with the plane. Every point on the sphere is mapped
to a unique point in the plane except of the north pole.
See \figref{stereo1} for an illustration.



In order to gain intuition, we first consider the one dimensional version of stereographic
projection between the unit circle $\Sp^1$ and and the real line. An example is shown
in \figref{stereo1D}.
Consider a line from the north pole of the unit circle $(0,1)$ that intersects
the unit circle at the point $(u,v)$ parameterized by $p(t)=(1-t)(0,1)+t(u,v).$
By considering the $y$-coordinate we determine the $t$ value were $p(t)$
intersects the $x$-axis, namely $t=\frac{1}{1-v}.$ Then, 
the $x$-coordinate where $p(t)$ intersects the real line is 
$$x=\frac{u}{1-v}.$$

\begin{figure}[htb]
	\centering
	\includegraphics[width=.3\textwidth]{curvature/stereo-1D}
	\caption{One dimensional stereographic projection.}
	\label{fig:stereo1D}
\end{figure}

A similar argument can be used to find the inverse map from
the real line to the unit circle. Given an $x$-value on the real line we consider
the parameterized line $p^{-1}(t)=(1-t)(0,1)+t(x,0)$ and compute
the $t$ value where $p^{-1}$ is on the unit circle and $t\neq 0.$ We find
$t=\frac{2}{1+x^2}$ and the point $(u,v)$ on the unit circle in terms of $x$ is
$$u=\frac{2x}{1+x^2}, v=\frac{x^2-1}{1+x^2}.$$

\begin{figure}[htb]
	\centering
	\includegraphics[width=.3\textwidth]{curvature/stereo}
	\caption{A point on the sphere is mapped to a point on the plane by stereographic projection.}
	\label{fig:stereo1}
\end{figure}


We can use this idea to find formulas for stereographic projection from
the unit sphere to the plane \cite{christian-notes}.
Consider the two sphere with the north pole $(0,0,1)$.
The line from the north pole $(0,0,1)$ that intersects $(x,y,z)\in \Sp^2$ parametrized by 
$p(t)=(1-t)(0,0,1)+t(x,y,z)$. By considering the $z$ coordinate we determine the $t$ value where this line
intersects $\R^2$, namely $t=\frac{1}{1-z}.$
This gives the desired map  and in equation form we have
$$p(x,y,z)\to \left(\frac{x}{1-z},\frac{y}{1-z}\right).$$

Many of us have more experience doing computations on $\R$ and $\R^2$ as opposed
to other objects. Stereographic projection enables us to transform many problems on
the circle and the sphere into problems in $\R$ and $\R^2$ respectively.
For example, suppose we wish to compute the integral of a constant function
$c(x,y)=1$ over the unit circle. 

	
